\chapter{Atividade 4 -- Integral pelo método do trapézio e análise de paralelização}

\label{chap:atvd04}

\section{Enunciado da tarefa}

\label{sec:atvd04-enunciado}

\begin{quote}

Implemente uma versão sequencial de um programa para o cálculo de integrais
definidas utilizando o método do trapézio.

O programa deve permitir calcular numericamente a integral de uma função
\(f(x)\) em um intervalo \([a,b]\), utilizando \(n\) subdivisões.

Em seguida, analise o código e discuta quais partes podem ser paralelizadas.
Também devem ser identificados possíveis desafios ou problemas relacionados
à paralelização desse problema, como dependências entre operações, acesso
concorrente a variáveis compartilhadas, necessidade de sincronização e
redução dos resultados parciais.

\end{quote}


\section{Objetivos e requisitos}

\label{sec:atvd04-objetivos}

A atividade teve como objetivo implementar uma versão sequencial do método
composto do trapézio e utilizar essa implementação como base para analisar
oportunidades e limitações de uma futura paralelização.

A parte prática foi deliberadamente mantida sequencial. A discussão de
paralelismo foi realizada a partir da estrutura do algoritmo e das
dependências observadas no laço responsável pela acumulação das contribuições
dos pontos internos.

Os principais objetivos considerados foram:

\begin{itemize}
    \item implementar corretamente o método composto do trapézio em C;
    \item utilizar aritmética inteira de 64 bits para tornar o cálculo
    determinístico e evitar efeitos de arredondamento em ponto flutuante;
    \item variar o número de subdivisões \(N\) e observar o comportamento do
    erro numérico e do tempo de execução;
    \item utilizar uma metodologia de temporização capaz de medir somente a
    região computacional de interesse;
    \item executar aquecimentos e múltiplas medições para reduzir o efeito de
    interferências externas;
    \item validar analiticamente cada resultado antes de considerá-lo válido
    para análise de desempenho;
    \item identificar quais dados seriam privados ou compartilhados em uma
    implementação paralela;
    \item discutir a condição de corrida associada ao acumulador da soma;
    \item identificar a operação de redução como mecanismo natural para uma
    futura implementação paralela.
\end{itemize}


\section{Fundamentação conceitual}

\label{sec:atvd04-fundamentacao}

\subsection{Integração numérica}

Uma integral definida

\begin{equation}
    I = \int_a^b f(x)\,dx
\end{equation}

pode ser interpretada geometricamente como a área algébrica delimitada pela
função \(f(x)\) no intervalo \([a,b]\).

Quando a integral não é calculada diretamente por uma expressão analítica,
métodos numéricos podem aproximar seu valor utilizando amostras discretas da
função.

Nesta atividade foi utilizado o método composto do trapézio.

\subsection{Método composto do trapézio}

Se o intervalo \([a,b]\) for dividido em \(N\) partes de mesma largura,
define-se

\begin{equation}
    h = \frac{b-a}{N}.
\end{equation}

Os pontos da discretização são dados por

\begin{equation}
    x_i = a + ih,
    \qquad
    i = 0,1,\ldots,N.
\end{equation}

Em cada subintervalo, a curva é aproximada por um segmento de reta. A área
aproximada do trapézio correspondente ao intervalo
\([x_i,x_{i+1}]\) é

\begin{equation}
    A_i =
    \frac{h}{2}
    \left[
        f(x_i)+f(x_{i+1})
    \right].
\end{equation}

A forma composta pode ser escrita como

\begin{equation}
    I_N =
    h
    \left[
        \frac{f(a)+f(b)}{2}
        +
        \sum_{i=1}^{N-1} f(a+ih)
    \right].
\end{equation}

Para a implementação inteira adotada nesta atividade, utilizou-se a forma
algebricamente equivalente

\begin{equation}
    I_N =
    h
    \sum_{i=1}^{N-1} f(x_i)
    +
    \frac{
        h\left[f(a)+f(b)\right]
    }{2}.
    \label{eq:atvd04-trapezio}
\end{equation}

Essa reorganização evita a formação intermediária de
\(2\sum f(x_i)\), reduzindo o risco de aproximação desnecessária do limite
representável por \texttt{int64\_t}.

\subsection{Função utilizada}

Foi adotada a função

\begin{equation}
    f(x)=x^2
\end{equation}

no intervalo

\begin{equation}
    [a,b]=[0,2\,400\,000].
\end{equation}

A integral analítica é

\begin{equation}
    \int_0^{2\,400\,000}x^2\,dx
    =
    \frac{(2\,400\,000)^3}{3},
\end{equation}

resultando em

\begin{equation}
    I_{\text{exato}}
    =
    4\,608\,000\,000\,000\,000\,000.
    \label{eq:atvd04-integral-exata}
\end{equation}

Esse valor permanece dentro do intervalo representável por um inteiro
assinado de 64 bits.

\subsection{Erro do método}

Para uma função suficientemente suave, o erro global do método composto do
trapézio é de ordem

\begin{equation}
    O(h^2).
\end{equation}

Para \(f(x)=x^2\), tem-se

\begin{equation}
    f''(x)=2,
\end{equation}

e o erro pode ser determinado exatamente por

\begin{equation}
    E_N =
    \frac{(b-a)h^2}{6}.
    \label{eq:atvd04-erro}
\end{equation}

Como \(x^2\) é convexa, o método do trapézio produz uma superestimação e,
portanto,

\begin{equation}
    I_N =
    I_{\text{exato}} + E_N.
    \label{eq:atvd04-resultado-esperado}
\end{equation}

Como

\begin{equation}
    h=\frac{b-a}{N},
\end{equation}

segue que

\begin{equation}
    E_N = O\left(\frac{1}{N^2}\right).
\end{equation}

Esse comportamento fornece também um critério independente de validação dos
resultados experimentais.

\subsection{Complexidade computacional}

O laço principal avalia todos os \(N-1\) pontos internos. Dessa forma, a
complexidade temporal é

\begin{equation}
    T(N)=O(N).
\end{equation}

A implementação não armazena os pontos da discretização nem os valores da
função em vetores. Cada ponto é produzido somente quando necessário.

Consequentemente, a memória adicional utilizada pelo kernel é constante:

\begin{equation}
    M(N)=O(1).
\end{equation}


\section{Interpretação didática da tarefa}

\label{sec:atvd04-interpretacao}

Embora o domínio matemático da atividade seja a integração numérica, o
objetivo didático relacionado à Computação de Alto Desempenho está na
estrutura da soma

\begin{equation}
    S =
    \sum_{i=1}^{N-1} f(a+ih).
\end{equation}

Cada ponto é obtido por

\begin{equation}
    x_i=a+ih.
\end{equation}

Portanto, o cálculo de \(x_i\) não depende de \(x_{i-1}\). Da mesma forma,
a avaliação

\begin{equation}
    f(x_i)
\end{equation}

não depende da avaliação realizada em outra iteração.

As contribuições individuais apresentam, assim, potencial para paralelismo
de dados.

Entretanto, a implementação sequencial contém a operação

\begin{verbatim}
soma += funcao_integrando(x_i);
\end{verbatim}

que estabelece uma dependência sobre o acumulador \texttt{soma}.

Em uma paralelização ingênua, múltiplas threads poderiam ler e modificar
essa variável simultaneamente, produzindo uma condição de corrida.

O problema pode ser representado conceitualmente como

\begin{figure}[htbp]
    \centering
    \begin{tikzpicture}[node distance=0.65cm and 0.7cm, font=\small]
        \node[thread, text width=2.7cm] (t1)
            {Thread 1\\\texttt{soma += valor\_1}};
        \node[thread, below=of t1, text width=2.7cm] (t2)
            {Thread 2\\\texttt{soma += valor\_2}};
        \node[thread, below=of t2, text width=2.7cm] (t3)
            {Thread 3\\\texttt{soma += valor\_3}};

        \node[componente, right=1.4cm of t2, text width=3.4cm] (soma)
            {Acumulador compartilhado\\\texttt{soma}};
        \node[decisao, below=0.8cm of soma, text width=2.7cm] (race)
            {Condição\\de corrida};

        \draw[fluxo] (t1.east) -- (soma.west);
        \draw[fluxo] (t2.east) -- (soma.west);
        \draw[fluxo] (t3.east) -- (soma.west);
        \draw[fluxo] (soma) -- (race);
    \end{tikzpicture}

    \caption{Atualização concorrente ingênua do acumulador compartilhado.}
    \label{fig:atvd04-corrida-soma}
\end{figure}

Uma estratégia mais apropriada consistiria em utilizar acumuladores privados

\begin{equation}
    S_1,S_2,\ldots,S_P
\end{equation}

e combiná-los posteriormente:

\begin{equation}
    S =
    \sum_{p=1}^{P}S_p.
\end{equation}

Esse padrão corresponde a uma operação de \emph{reduction}.

A atividade, portanto, estabelece a base conceitual para uma implementação
paralela posterior: o cálculo dos pontos é independente, enquanto a
acumulação exige tratamento específico.


\section{Modelagem da solução}

\label{sec:atvd04-modelagem}

\subsection{Estrutura geral}

A solução foi organizada em quatro etapas principais:

\begin{enumerate}
    \item definição do problema matemático e dos valores de \(N\);
    \item execução do kernel sequencial do método do trapézio;
    \item realização dos aquecimentos e medições temporais;
    \item validação matemática e apresentação dos resultados.
\end{enumerate}

A Figura~\ref{fig:atvd04-modelagem} resume essa modelagem em termos de
etapas e responsabilidades.

\begin{figure}[htbp]
    \centering
    \begin{tikzpicture}[node distance=0.65cm, font=\small]
        \node[processo, text width=5.0cm] (problema)
            {Definir problema\\\(f(x)=x^2\), \([a,b]\) e \(N\)};
        \node[componente, below=of problema, text width=5.0cm] (kernel)
            {Executar kernel sequencial\\\texttt{integral\_trapezio()}};
        \node[componente, below=of kernel, text width=5.0cm] (medicao)
            {Aquecimentos e medições\\\texttt{CLOCK\_MONOTONIC}};
        \node[componente, below=of medicao, text width=5.0cm] (validacao)
            {Validar resultado\\erro analítico e mediana};
        \node[componente, below=of validacao, text width=5.0cm] (saida)
            {Apresentar métricas\\tempo, erro e ns/trap};

        \draw[fluxo] (problema) -- (kernel);
        \draw[fluxo] (kernel) -- (medicao);
        \draw[fluxo] (medicao) -- (validacao);
        \draw[fluxo] (validacao) -- (saida);
    \end{tikzpicture}

    \caption{Modelagem UML da solução sequencial da Atividade 4.}
    \label{fig:atvd04-modelagem}
\end{figure}

Os valores de \(N\) utilizados foram

\begin{equation}
\begin{split}
N \in \{&
24\,000,\,
120\,000,\,
240\,000,\,
400\,000,\\
&
800\,000,\,
1\,200\,000,\,
2\,400\,000
\}.
\end{split}
\end{equation}

Todos esses valores dividem exatamente o comprimento do intervalo

\begin{equation}
    b-a=2\,400\,000,
\end{equation}

garantindo que

\begin{equation}
    h=\frac{b-a}{N}
\end{equation}

permaneça inteiro.

\subsection{Principais funções e componentes}

A implementação foi decomposta nas seguintes responsabilidades:

\begin{description}

    \item[\texttt{funcao\_integrando()}]
    Calcula \(f(x)=x^2\).

    \item[\texttt{integral\_trapezio()}]
    Implementa o kernel sequencial do método composto do trapézio.

    \item[\texttt{obter\_tempo()}]
    Obtém um instante utilizando \texttt{CLOCK\_MONOTONIC}.

    \item[\texttt{diferenca\_ns()}]
    Calcula o intervalo entre dois instantes em nanossegundos.

    \item[\texttt{comparar\_uint64()}]
    Implementa a função auxiliar utilizada na ordenação dos tempos.

    \item[\texttt{calcular\_mediana()}]
    Calcula a mediana das vinte medições de uma configuração.

    \item[\texttt{resultado\_esperado()}]
    Calcula analiticamente o resultado que o método do trapézio deve
    produzir para determinado \(N\).

    \item[\texttt{executar\_benchmark()}]
    Coordena aquecimentos, medições, validação, normalização e impressão
    dos resultados.

    \item[\texttt{main()}]
    Percorre os diferentes valores de \(N\) e coordena o experimento.

\end{description}


\section{Decisões de projeto e trade-offs}

\label{sec:atvd04-decisoes}

\subsection{Uso de \texttt{int64\_t}}

Foi utilizada aritmética inteira em vez de \texttt{float} ou \texttt{double}
no kernel.

Essa decisão teve como objetivos:

\begin{itemize}
    \item evitar diferenças decorrentes de arredondamento em ponto
    flutuante;
    \item tornar a validação determinística;
    \item facilitar a futura análise de redução;
    \item permitir comparação exata com o resultado teórico do método.
\end{itemize}

O uso de \texttt{int} de 32 bits não seria suficiente. Por exemplo,

\begin{equation}
    (2\,400\,000)^2
    =
    5\,760\,000\,000\,000,
\end{equation}

valor já superior ao limite de um inteiro assinado de 32 bits.

Por esse motivo foi adotado \texttt{int64\_t}.

\subsection{Escolha do intervalo}

O limite superior \(2\,400\,000\) foi escolhido de forma que:

\begin{itemize}
    \item existissem diversos divisores adequados para \(N\);
    \item \(h\) permanecesse inteiro;
    \item o valor da integral ainda coubesse em \texttt{int64\_t}.
\end{itemize}

Essa escolha também impôs uma limitação: para manter \(h\geq1\), o maior
valor possível de \(N\) é \(2\,400\,000\).

Essa restrição contribuiu para a adoção posterior de repetições internas
durante a temporização.

\subsection{Cálculo direto de \(x_i\)}

Os pontos são calculados por

\begin{equation}
    x_i=a+ih
\end{equation}

em vez de utilizar a recorrência

\begin{equation}
    x_{i+1}=x_i+h.
\end{equation}

Essa decisão evita introduzir uma dependência artificial entre iterações e
torna explícito que a dependência principal do laço está apenas na
acumulação da soma.

\subsection{Ausência de vetores intermediários}

Não foram armazenados vetores contendo os pontos \(x_i\) ou os valores
\(f(x_i)\).

Essa decisão mantém

\begin{equation}
    M(N)=O(1)
\end{equation}

e evita introduzir efeitos de alocação e acesso à memória que não fazem
parte do objetivo principal da atividade.

\subsection{Nível de otimização}

A compilação experimental foi realizada com otimização

\begin{verbatim}
-O2
\end{verbatim}

em vez de \texttt{-O0}.

O uso de \texttt{-O0} produziria um executável artificialmente pouco
otimizado e não representativo de uma execução sequencial realista.

Entretanto, foi utilizada a opção

\begin{verbatim}
-fno-tree-vectorize
\end{verbatim}

para impedir que o compilador transformasse automaticamente o laço escalar
em código SIMD.

Dessa forma, a execução permaneceu sequencial e escalar para fins da
atividade, ainda permitindo otimizações convencionais do compilador.

\subsection{Proteção contra otimizações interprocedurais}

Durante o desenvolvimento, uma versão inicial do benchmark apresentou
tempos iguais a zero, mesmo para os maiores valores de \(N\).

A análise indicou que impedir apenas o \emph{inlining} não era suficiente.
O compilador podia reconhecer que chamadas consecutivas do kernel,
realizadas com os mesmos argumentos, produziriam sempre o mesmo resultado
e reutilizar a computação.

Por esse motivo, o kernel foi marcado, no GCC, com

\begin{verbatim}
__attribute__((noinline, noipa))
\end{verbatim}

O atributo \texttt{noinline} impede a incorporação da função ao chamador,
enquanto \texttt{noipa} impede otimizações interprocedurais capazes de
eliminar ou reutilizar chamadas repetidas.

Adicionalmente, o resultado é armazenado em uma variável
\texttt{volatile} após a região cronometrada, tornando seu uso observável.

O acumulador principal, entretanto, permanece não-\texttt{volatile}, de
forma a não alterar artificialmente o desempenho do kernel.

\subsection{Repetições internas}

Após a correção das otimizações interprocedurais, as execuções individuais
ainda eram curtas, variando da ordem de microssegundos até frações de
milissegundo.

Para aumentar a duração efetivamente observada pelo relógio, cada medição
passou a executar

\begin{equation}
    R=1000
\end{equation}

chamadas consecutivas do kernel.

O tempo medido para o lote é posteriormente normalizado:

\begin{equation}
    T_{\text{execução}}
    =
    \frac{T_{\text{lote}}}{1000}.
\end{equation}

Essa estratégia reduz a influência relativa do overhead de temporização sem
alterar o algoritmo analisado.


\section{Representação estrutural da implementação}

\label{sec:atvd04-estrutura}

A Figura~\ref{fig:atvd04-estrutura} apresenta uma representação estrutural
simplificada da implementação.

\begin{figure}[htbp]
    \centering
    \resizebox{0.96\textwidth}{!}{%
    \begin{tikzpicture}[node distance=0.8cm and 0.7cm, font=\small]
        \node[processo, text width=3.4cm] (main)
            {\texttt{main()}\\valores de \(N\)\\configuração};

        \node[componente, right=of main, text width=3.8cm] (bench)
            {\texttt{executar\_benchmark()}\\coordena cada caso};

        \node[componente, below left=of bench, text width=3.4cm] (warmup)
            {3 aquecimentos\\descartados};
        \node[componente, below=of bench, text width=3.4cm] (medicoes)
            {20 medições\\1000 chamadas};
        \node[componente, below right=of bench, text width=3.4cm] (pos)
            {mediana, validação\\erro e ns/trap};

        \node[processo, below=of medicoes, text width=3.6cm] (kernel)
            {\texttt{integral\_trapezio()}\\kernel sequencial};
        \node[componente, right=of pos, text width=3.3cm] (saida)
            {impressão\\dos resultados};

        \draw[fluxo] (main) -- (bench);
        \draw[fluxo] (bench) -- (warmup);
        \draw[fluxo] (bench) -- (medicoes);
        \draw[fluxo] (bench) -- (pos);
        \draw[fluxo] (warmup) |- (kernel);
        \draw[fluxo] (medicoes) -- (kernel);
        \draw[fluxo] (pos) -- (saida);
    \end{tikzpicture}%
    }

    \caption{Estrutura geral da implementação da Atividade 4.}
    \label{fig:atvd04-estrutura}
\end{figure}


\section{Fluxo de execução}

\label{sec:atvd04-fluxo}

O fluxo experimental é apresentado na
Figura~\ref{fig:atvd04-fluxo}.

\begin{figure}[H]
    \centering
    \begin{tikzpicture}[node distance=0.48cm and 0.8cm, font=\small]
        \node[processo, text width=3.8cm] (inicio) {Início};
        \node[componente, below=of inicio, text width=4.3cm] (config)
            {Configuração experimental};
        \node[componente, below=of config, text width=4.3cm] (selecionar)
            {Selecionar valor de \(N\)};
        \node[componente, below=of selecionar, text width=4.3cm] (aquec)
            {Executar 3 aquecimentos};
        \node[componente, below=of aquec, text width=4.3cm] (medir)
            {20 medições válidas\\\(t_0\), 1000 kernels, \(t_1\)};
        \node[componente, below=of medir, text width=4.3cm] (mediana)
            {Calcular mediana\\e normalizar tempo};
        \node[decisao, below=of mediana, text width=2.5cm] (validar)
            {Resultado\\correto?};
        \node[componente, below left=0.55cm and 1.45cm of validar, text width=3.3cm] (metricas)
            {Calcular métricas\\e imprimir};
        \node[componente, below right=0.55cm and 1.45cm of validar, text width=3.0cm] (abortar)
            {Abortar};
        \node[decisao, below=1.6cm of metricas, text width=2.6cm] (outro)
            {Existe\\outro \(N\)?};
        \node[processo, below=of outro, text width=3.4cm] (fim)
            {Fim};

        \draw[fluxo] (inicio) -- (config);
        \draw[fluxo] (config) -- (selecionar);
        \draw[fluxo] (selecionar) -- (aquec);
        \draw[fluxo] (aquec) -- (medir);
        \draw[fluxo] (medir) -- (mediana);
        \draw[fluxo] (mediana) -- (validar);
        \draw[fluxo] (validar) -| node[pos=0.25, above] {sim} (metricas);
        \draw[fluxo] (validar) -| node[pos=0.25, above] {não} (abortar);
        \draw[fluxo] (metricas) -- (outro);
        \draw[fluxo] (outro) -- node[right] {não} (fim);
        \draw[fluxo] (outro.west) -- ++(-0.8,0) |- (selecionar.west);
        \node[above] at ($(outro.west)+(-0.15,0)$) {sim};
    \end{tikzpicture}

    \caption{Fluxo de execução do benchmark sequencial.}
    \label{fig:atvd04-fluxo}
\end{figure}


\section{Metodologia experimental}

\label{sec:atvd04-metodologia}

\subsection{Ambiente de execução}

Os experimentos da Atividade 4 foram executados em ambiente Windows,
utilizando terminal MSYS2/UCRT64 e compilador GCC.

A compilação utilizada foi equivalente a:

\begin{verbatim}
gcc -std=c11 -O2 -Wall -Wextra -Wpedantic \
    -fno-tree-vectorize tarefa4.c -o tarefa4.exe
\end{verbatim}

Não foi registrado, durante a execução experimental consolidada, o número
de versão do compilador no conjunto de resultados utilizado neste capítulo;
por isso esse dado não é atribuído posteriormente.

\subsection{Configuração das medições}

Para cada valor de \(N\) foram executados:

\begin{equation}
    3
\end{equation}

aquecimentos, descartados das estatísticas, seguidos de

\begin{equation}
    20
\end{equation}

medições válidas.

Cada medição executou internamente

\begin{equation}
    1000
\end{equation}

chamadas do kernel.

O tempo do lote foi obtido por

\begin{verbatim}
clock_gettime(CLOCK_MONOTONIC, ...)
\end{verbatim}

e normalizado para uma única execução.

A estatística representativa utilizada foi a mediana das vinte medições.

A escolha da mediana reduz a influência de observações isoladas afetadas
por interrupções, escalonamento do sistema operacional ou processos em
segundo plano.

\subsection{Região cronometrada}

Dentro da região medida permaneceram somente:

\begin{itemize}
    \item cálculo da integral;
    \item cálculo de \(h\);
    \item cálculo dos pontos internos;
    \item avaliação de \(f(x_i)\);
    \item acumulação da soma;
    \item combinação com os extremos.
\end{itemize}

Operações como impressão, validação, ordenação dos tempos, cálculo da
mediana e formatação do resultado permaneceram fora da região cronometrada.

\subsection{Métricas}

Foram utilizadas três grandezas principais:

\paragraph{Tempo mediano}

\begin{equation}
    T_{\text{med}}
\end{equation}

correspondente ao tempo normalizado de uma execução.

\paragraph{Erro absoluto}

\begin{equation}
    E_{\text{abs}}
    =
    \left|
        I_N-I_{\text{exato}}
    \right|.
\end{equation}

\paragraph{Tempo normalizado por trapézio}

\begin{equation}
    C_N =
    \frac{T_{\text{med}}}{N}.
\end{equation}

A métrica foi apresentada em nanossegundos por trapézio:

\begin{equation}
    \text{ns/trap}.
\end{equation}

Como o algoritmo possui complexidade \(O(N)\), espera-se que essa razão
permaneça aproximadamente constante para execuções suficientemente
representativas.


\section{Código-fonte desenvolvido}

\label{sec:atvd04-codigo}

\lstinputlisting[
    style=codigoC,
    caption={Código-fonte da Atividade 4.},
    label={lst:atvd04-codigo}
]{../ATVD4/tarefa4.c}


\section{Validação da implementação}

\label{sec:atvd04-validacao}

A validação foi realizada utilizando simultaneamente o valor analítico da
integral e a expressão exata do erro do método para \(f(x)=x^2\).

Para cada \(N\), calculou-se

\begin{equation}
    h=\frac{2\,400\,000}{N}
\end{equation}

e

\begin{equation}
    E_N=
    \frac{2\,400\,000\,h^2}{6}.
\end{equation}

A Tabela~\ref{tab:atvd04-validacao} apresenta os valores esperados.

\begin{table}[htbp]
    \centering
    \caption{Validação analítica do erro do método do trapézio.}
    \label{tab:atvd04-validacao}
    \begin{tabular}{r|r|r}
        \hline
        \(N\) & \(h\) & Erro esperado \\
        \hline
        24\,000    & 100 & 4\,000\,000\,000 \\
        120\,000   & 20  & 160\,000\,000 \\
        240\,000   & 10  & 40\,000\,000 \\
        400\,000   & 6   & 14\,400\,000 \\
        800\,000   & 3   & 3\,600\,000 \\
        1\,200\,000 & 2  & 1\,600\,000 \\
        2\,400\,000 & 1  & 400\,000 \\
        \hline
    \end{tabular}
\end{table}

Todos os valores calculados experimentalmente coincidiram exatamente com
esses valores esperados.

A validação, portanto, não utilizou tolerância de ponto flutuante. A
comparação foi realizada diretamente com valores inteiros.

Esse resultado confirmou simultaneamente:

\begin{itemize}
    \item a correta implementação do método;
    \item a ausência de overflow nos casos utilizados;
    \item a validade dos valores escolhidos para \(N\);
    \item a convergência esperada do método.
\end{itemize}


\section{Resultados experimentais}

\label{sec:atvd04-resultados}

A Tabela~\ref{tab:atvd04-resultados} apresenta os resultados finais obtidos
após a introdução das mil repetições internas por medição.

\begin{table}[htbp]
    \centering
    \caption{Resultados experimentais finais da Atividade 4.}
    \label{tab:atvd04-resultados}
    \small
    \begin{tabular}{r|r|r|r|r}
        \hline
        \(N\) &
        Integral &
        Erro absoluto &
        Mediana (ms) &
        ns/trap \\
        \hline
        24\,000 &
        4608000004000000000 &
        4000000000 &
        0.005066 &
        0.211 \\

        120\,000 &
        4608000000160000000 &
        160000000 &
        0.025484 &
        0.212 \\

        240\,000 &
        4608000000040000000 &
        40000000 &
        0.050818 &
        0.212 \\

        400\,000 &
        4608000000014400000 &
        14400000 &
        0.084860 &
        0.212 \\

        800\,000 &
        4608000000003600000 &
        3600000 &
        0.171002 &
        0.214 \\

        1\,200\,000 &
        4608000000001600000 &
        1600000 &
        0.254809 &
        0.212 \\

        2\,400\,000 &
        4608000000000400000 &
        400000 &
        0.511678 &
        0.213 \\
        \hline
    \end{tabular}
\end{table}

A integral converge progressivamente para

\begin{equation}
    4\,608\,000\,000\,000\,000\,000,
\end{equation}

enquanto o erro absoluto diminui conforme o número de subdivisões aumenta.


\section{Visualização dos resultados}

\label{sec:atvd04-visualizacao}

Os resultados apresentam duas tendências distintas e complementares.

\begin{center}
\fbox{
\begin{minipage}{0.84\textwidth}
\centering
\textbf{Comportamento numérico}

\vspace{0.2cm}

\(N \uparrow\)
\(\quad\Longrightarrow\quad\)
\(h \downarrow\)
\(\quad\Longrightarrow\quad\)
\(E_{\text{abs}}\downarrow\)

\vspace{0.2cm}

\[
E_{\text{abs}}
=
O(N^{-2})
\]

\vspace{0.35cm}

\textbf{Comportamento computacional}

\vspace{0.2cm}

\(N \uparrow\)
\(\quad\Longrightarrow\quad\)
número de iterações \(\uparrow\)
\(\quad\Longrightarrow\quad\)
tempo \(\uparrow\)

\vspace{0.2cm}

\[
T(N)=O(N)
\]

\vspace{0.2cm}

\[
\frac{T(N)}{N}
\approx
0.21\ \text{ns/trap}
\]
\end{minipage}
}
\end{center}

A estabilidade da coluna \texttt{ns/trap} constitui uma forma normalizada
de visualizar o comportamento linear do kernel.

Os valores obtidos foram

\begin{equation}
    0.211,\,
    0.212,\,
    0.212,\,
    0.212,\,
    0.214,\,
    0.212,\,
    0.213
    \ \text{ns/trap}.
\end{equation}

A variação é pequena ao longo de toda a faixa experimental.


\section{Análise e discussão}

\label{sec:atvd04-discussao}

\subsection{Comportamento do erro}

Os resultados confirmam a ordem de convergência esperada.

Quando \(N\) passa de

\begin{equation}
    120\,000
\end{equation}

para

\begin{equation}
    240\,000,
\end{equation}

o número de subdivisões dobra.

O erro, entretanto, passa de

\begin{equation}
    160\,000\,000
\end{equation}

para

\begin{equation}
    40\,000\,000,
\end{equation}

ou seja, é reduzido por um fator de quatro.

Isso é compatível com

\begin{equation}
    E_N=O(N^{-2}).
\end{equation}

A mesma propriedade pode ser observada nos demais pares de configurações
compatíveis.

\subsection{Comportamento temporal}

O tempo de execução apresentou crescimento aproximadamente proporcional ao
número de subdivisões.

Por exemplo, entre

\begin{equation}
    N=120\,000
\end{equation}

e

\begin{equation}
    N=240\,000,
\end{equation}

o número de subdivisões dobra, enquanto o tempo passa de

\begin{equation}
    0.025484\ \text{ms}
\end{equation}

para

\begin{equation}
    0.050818\ \text{ms}.
\end{equation}

A razão entre os tempos é aproximadamente

\begin{equation}
    \frac{0.050818}{0.025484}
    \approx
    1.994.
\end{equation}

Analogamente, entre \(400\,000\) e \(800\,000\),

\begin{equation}
    \frac{0.171002}{0.084860}
    \approx
    2.015,
\end{equation}

e entre \(1\,200\,000\) e \(2\,400\,000\),

\begin{equation}
    \frac{0.511678}{0.254809}
    \approx
    2.008.
\end{equation}

Esses resultados são compatíveis com a complexidade

\begin{equation}
    T(N)=O(N).
\end{equation}

\subsection{Custo normalizado}

O custo por trapézio permaneceu aproximadamente constante em torno de

\begin{equation}
    0.21\ \text{ns/trap}.
\end{equation}

Esse resultado reforça a interpretação de crescimento linear.

O pequeno intervalo observado,

\begin{equation}
    0.211
    \leq
    C_N
    \leq
    0.214,
\end{equation}

é compatível com variações normais de execução do sistema.

\subsection{Efeito das repetições internas}

As execuções individuais do kernel apresentaram tempos curtos. No menor
caso, o tempo normalizado foi

\begin{equation}
    0.005066\ \text{ms},
\end{equation}

enquanto no maior foi

\begin{equation}
    0.511678\ \text{ms}.
\end{equation}

Como cada medição executou mil chamadas, as regiões efetivamente observadas
pelo relógio apresentaram duração mediana aproximada entre

\begin{equation}
    5.066\ \text{ms}
\end{equation}

e

\begin{equation}
    511.678\ \text{ms}.
\end{equation}

Isso reduziu significativamente a influência relativa do custo das chamadas
ao temporizador.

Além disso, os tempos normalizados permaneceram próximos aos valores
observados antes da introdução do batching, fornecendo evidência de que a
alteração aumentou a robustez da medição sem modificar a tendência do
kernel.

\subsection{Análise de paralelização}

O laço principal pode ser representado conceitualmente por

\begin{verbatim}
for i = 1 .. N-1
    x_i   = a + i*h
    valor = f(x_i)
    soma  = soma + valor
\end{verbatim}

As duas primeiras operações podem ser executadas independentemente para
diferentes valores de \(i\).

O problema aparece na terceira operação.

A atualização

\begin{equation}
    S \leftarrow S + v_i
\end{equation}

envolve leitura, operação aritmética e escrita.

Se diferentes threads realizarem essas etapas simultaneamente sobre o mesmo
acumulador, atualizações podem ser perdidas.

A versão paralela ingênua teria, portanto, uma condição de corrida.

Uma possível decomposição seria:

\begin{figure}[htbp]
    \centering
    \begin{tikzpicture}[node distance=0.75cm and 0.8cm, font=\small]
        \node[processo, text width=4.8cm] (iter)
            {Iterações independentes\\\(i=1,\ldots,N-1\)};

        \node[thread, below left=of iter, text width=2.6cm] (t1)
            {Thread 1\\subfaixa 1};
        \node[thread, below=of iter, text width=2.6cm] (t2)
            {Thread 2\\subfaixa 2};
        \node[thread, below right=of iter, text width=2.6cm] (tp)
            {Thread \(P\)\\subfaixa \(P\)};

        \node[componente, below=of t1, text width=2.4cm] (s1) {\(S_1\)};
        \node[componente, below=of t2, text width=2.4cm] (s2) {\(S_2\)};
        \node[componente, below=of tp, text width=2.4cm] (sp) {\(S_P\)};

        \node[processo, below=1.0cm of s2, text width=4.0cm] (red)
            {Redução\\\(S=\sum_{p=1}^{P}S_p\)};
        \node[componente, below=of red, text width=3.0cm] (s)
            {Soma final\\\(S\)};

        \draw[fluxo] (iter) -- (t1);
        \draw[fluxo] (iter) -- (t2);
        \draw[fluxo] (iter) -- (tp);
        \draw[fluxo] (t1) -- (s1);
        \draw[fluxo] (t2) -- (s2);
        \draw[fluxo] (tp) -- (sp);
        \draw[fluxo] (s1) |- (red);
        \draw[fluxo] (s2) -- (red);
        \draw[fluxo] (sp) |- (red);
        \draw[fluxo] (red) -- (s);
    \end{tikzpicture}

    \caption{Decomposição conceitual do acumulador por redução.}
    \label{fig:atvd04-reduction}
\end{figure}

Nesse modelo, cada thread mantém um acumulador privado durante sua parcela
do trabalho. Somente posteriormente os resultados parciais são combinados.

Essa abordagem evita a atualização concorrente de uma única variável a cada
iteração.

\subsection{Classificação das variáveis}

Em uma eventual versão paralela, as variáveis poderiam ser classificadas
como apresentado na Tabela~\ref{tab:atvd04-variaveis}.

\begin{table}[htbp]
    \centering
    \caption{Classificação conceitual das variáveis em uma paralelização.}
    \label{tab:atvd04-variaveis}

    \begin{tabular}{l|l}
        \hline
        Elemento & Classificação conceitual \\
        \hline
        \(a\) & compartilhado, somente leitura \\
        \(b\) & compartilhado, somente leitura \\
        \(N\) & compartilhado, somente leitura \\
        \(h\) & compartilhado, somente leitura \\
        \(i\) & privado por thread \\
        \(x_i\) & privado por thread \\
        \(f(x_i)\) & privado por thread \\
        \texttt{soma} & variável de redução \\
        \hline
    \end{tabular}
\end{table}

Essa classificação mostra que a maior parte do estado do algoritmo não
exige sincronização. O ponto crítico é a combinação dos resultados
parciais.


\section{Limitações}

\label{sec:atvd04-limitacoes}

Os resultados devem ser interpretados considerando algumas limitações.

Primeiramente, o experimento utiliza uma função específica,

\begin{equation}
    f(x)=x^2,
\end{equation}

cujo custo por avaliação é baixo e aproximadamente uniforme. Funções mais
complexas podem alterar significativamente a relação entre custo aritmético
e overhead de paralelização.

A adoção de aritmética inteira melhora a reprodutibilidade do experimento,
mas não representa todas as aplicações de integração numérica, que
frequentemente utilizam ponto flutuante.

A opção \texttt{-fno-tree-vectorize} impede vetorização automática. Essa
decisão é adequada para manter um baseline escalar, porém significa que os
tempos não representam necessariamente o máximo desempenho sequencial que
o compilador poderia produzir em uma aplicação real.

Os atributos \texttt{noinline} e \texttt{noipa} também constituem uma
intervenção metodológica destinada a preservar a validade do benchmark.
Eles evitam que o compilador elimine trabalho que precisa ser medido, mas
restringem algumas otimizações possíveis em um programa não experimental.

As mil repetições internas utilizam sempre os mesmos parâmetros. Essa
estratégia aumenta a duração da região cronometrada e reduz o impacto do
temporizador, porém não representa uma sequência de problemas independentes
com dados diferentes.

Os experimentos foram executados em um sistema operacional de propósito
geral. Escalonamento de processos, frequência dinâmica do processador,
interrupções e outros serviços podem produzir pequenas variações entre
medições.

Por fim, a atividade não implementa efetivamente uma versão paralela.
Assim, as discussões sobre condição de corrida, sincronização e redução são
derivadas da estrutura do algoritmo, conforme solicitado pelo enunciado,
e não de medições paralelas nesta atividade.


\section{Conclusões da atividade}

\label{sec:atvd04-conclusao}

A implementação sequencial do método composto do trapézio foi validada
analiticamente para todos os valores de \(N\) avaliados.

O uso de \(f(x)=x^2\), aritmética \texttt{int64\_t} e subdivisões que
mantêm \(h\) inteiro permitiu comparar exatamente o resultado calculado com
o valor teoricamente esperado.

O erro absoluto diminuiu conforme

\begin{equation}
    O(N^{-2}),
\end{equation}

em concordância com a ordem de erro do método composto do trapézio.

Em relação ao custo computacional, os tempos experimentais apresentaram
crescimento aproximadamente linear com \(N\), enquanto o custo normalizado
permaneceu entre aproximadamente \(0.211\) e \(0.214\) ns por trapézio.

A atividade também permitiu identificar a estrutura de paralelismo do
problema. O cálculo de cada ponto

\begin{equation}
    x_i=a+ih
\end{equation}

e sua respectiva avaliação são independentes. Entretanto, a atualização do
acumulador representa uma dependência de redução.

Uma paralelização direta utilizando um único acumulador compartilhado pode
produzir condição de corrida. A estratégia conceitualmente mais adequada é
a utilização de acumuladores privados seguida de uma etapa de redução.

Dessa forma, a atividade estabelece tanto um baseline sequencial validado
quanto a fundamentação necessária para uma implementação paralela posterior.


\section{Síntese}

\label{sec:atvd04-defesa}

Para uma apresentação oral, os principais pontos da atividade podem ser
sintetizados da seguinte forma:

\begin{enumerate}
    \item A tarefa exigiu implementar somente a versão sequencial e discutir
    teoricamente sua paralelização.

    \item Foi utilizada a função
    \[
        f(x)=x^2
    \]
    no intervalo
    \[
        [0,2\,400\,000].
    \]

    \item O tipo \texttt{int64\_t} foi adotado para evitar efeitos de
    arredondamento e permitir validação inteira exata.

    \item O método implementado foi o trapézio composto, com
    \[
        x_i=a+ih.
    \]
    Essa forma evita dependência entre pontos consecutivos.

    \item Foram utilizados sete valores de \(N\), todos divisores de
    \(2\,400\,000\), garantindo passo \(h\) inteiro.

    \item O valor exato da integral é
    \[
        4\,608\,000\,000\,000\,000\,000.
    \]

    \item Para \(f(x)=x^2\), o erro pode ser calculado exatamente por
    \[
        E_N=\frac{(b-a)h^2}{6}.
    \]
    Todos os resultados experimentais coincidiram com esse valor.

    \item Foram realizados 3 aquecimentos e 20 medições para cada \(N\).

    \item Cada medição executou 1000 chamadas do kernel. O tempo do lote foi
    dividido por 1000 para obter o tempo equivalente de uma execução.

    \item Foi utilizado
    \texttt{clock\_gettime(CLOCK\_MONOTONIC)} para medir tempo decorrido.

    \item Durante o desenvolvimento, uma primeira versão apresentou tempos
    zerados devido a otimizações interprocedurais. O problema foi corrigido
    utilizando \texttt{noinline}, \texttt{noipa} e consumo observável do
    resultado.

    \item A compilação utilizou \texttt{-O2}, mas a vetorização automática
    foi desabilitada com \texttt{-fno-tree-vectorize} para preservar um
    baseline escalar.

    \item O erro diminuiu com comportamento
    \[
        O(N^{-2}).
    \]

    \item O tempo aumentou aproximadamente de forma linear:
    \[
        T(N)=O(N).
    \]

    \item O custo normalizado permaneceu próximo de
    \[
        0.21\ \text{ns/trap}.
    \]

    \item O cálculo de cada \(f(x_i)\) pode ser paralelizado, pois as
    iterações são independentes.

    \item A variável \texttt{soma} não pode ser atualizada ingenuamente por
    várias threads, pois isso produziria uma condição de corrida.

    \item A estratégia conceitualmente mais adequada é utilizar somas
    privadas por thread e combinar os resultados posteriormente por uma
    operação de redução.
\end{enumerate}
